\chapter{Μέθοδοι Πρόβλεψης}
\label{ch:methods}

\section{Γραμμές Βάσης (Baselines)}
\label{sec:baselines}

Τρεις κλασικές γραμμές βάσης χρησιμοποιούνται για αξιολόγηση:

\begin{description}
    \item[\textbf{Naive-1:}] $\hat{y}_{t+1} = y_t$. Πρόβλεψη ίση με την τελευταία παρατηρηθείσα τιμή.
    \item[\textbf{Naive-24:}] $\hat{y}_{t+1} = y_{t-23}$. Πρόβλεψη ίση με την τιμή της ίδιας ώρας
    χθες. Αποτελεί φυσικό ανταγωνιστή για MIMO H=24 (προβλέπει ολόκληρη ημέρα από χθες).
    \item[\textbf{Seasonal Profile:}] Η πρόβλεψη βασίζεται στον μέσο όρο ιστορικών τιμών για την
    ίδια ώρα και ημέρα εβδομάδας.
\end{description}

Τα Naive-1/24 είναι ιδιαίτερα χρήσιμα γιατί σε ορισμένες περιόδους υψηλής μεταβλητότητας
(π.χ.\ Q4 2025) ξεπερνούν τα ML μοντέλα σε εβδομαδιαίο ορίζοντα.


\section{Αλγόριθμοι Μηχανικής Μάθησης}
\label{sec:ml_algorithms}

\subsection{LightGBM}

Το LightGBM~\cite{ke2017lgbm} αποτελεί gradient boosting framework με leaf-wise ανάπτυξη δέντρων.
Στην παρούσα εργασία, ο βασικός configuration που χρησιμοποιείται είναι:

\begin{table}[H]
\centering
\caption{Παράμετροι LightGBM (standard CL models).}
\label{tab:lgbm_params}
\begin{tabular}{ll}
\toprule
\textbf{Παράμετρος} & \textbf{Τιμή} \\
\midrule
\texttt{n\_estimators} & 1000 \\
\texttt{learning\_rate} & 0.05 (default) \\
\texttt{num\_leaves} & 31 (default) \\
\texttt{objective} & \texttt{regression} (MAE) \\
\texttt{early\_stopping\_rounds} & 50 \\
\bottomrule
\end{tabular}
\end{table}

Τα βελτιστοποιημένα μοντέλα (Optuna) χρησιμοποιούν:
\texttt{n\_est=900, lr=0.027, num\_leaves=26, min\_child=54, subsample=0.925, colsample=0.641}.

\subsection{XGBoost}

Το XGBoost~\cite{chen2016xgboost} χρησιμοποιεί level-wise ανάπτυξη δέντρων με second-order
gradient approximation. Τυπικές παράμετροι:

\begin{table}[H]
\centering
\caption{Παράμετροι XGBoost (standard CL models).}
\label{tab:xgb_params}
\begin{tabular}{ll}
\toprule
\textbf{Παράμετρος} & \textbf{Τιμή} \\
\midrule
\texttt{n\_estimators} & 1000 \\
\texttt{learning\_rate} & 0.05 \\
\texttt{max\_depth} & 6 \\
\texttt{objective} & \texttt{reg:absoluteerror} \\
\texttt{early\_stopping\_rounds} & 50 \\
\bottomrule
\end{tabular}
\end{table}

\subsection{Random Forest}

Το Random Forest~\cite{breiman2001rf} εκπαιδεύεται με:
\begin{itemize}
    \item \texttt{n\_estimators=400} (600 για MIMO two-stage)
    \item \texttt{max\_features="sqrt"} (περίπου 12 από τα 152 χαρακτηριστικά ανά split, κρίσιμο
    για ταχύτητα στα Windows χωρίς loky deadlock)
    \item \texttt{max\_depth=16, min\_samples\_split=4, n\_jobs=1}
\end{itemize}

\subsection{Support Vector Regression}

Η SVR~\cite{smola2004} χρησιμοποιεί RBF kernel με $C=10$, $\varepsilon=0.1$, προεπεξεργασία
\texttt{StandardScaler}. Λόγω του $O(n^2)$ κόστους εκπαίδευσης, τα MIMO SVR μοντέλα εκπαιδεύτηκαν
σε 5.000 τυχαία δείγματα (από τα ~76.000 διαθέσιμα), κάτι που επηρεάζει αρνητικά την επίδοσή τους.

\subsection{Multi-Layer Perceptron (MLP)}

Το MLP υλοποιείται με PyTorch και sklearn wrapper (\texttt{MLPRegressor}):
\begin{itemize}
    \item \textbf{CL/OL:} hidden=(512,256), lr=5e-4, max\_iter=300, early stopping
    \item \textbf{MIMO:} hidden=(256,128), lr=5e-4, max\_iter=400, \texttt{StandardScaler} προεπεξεργασία
    \item \textbf{MLP-Optuna CL:} hidden=[64] τιμή, hidden=[128,64] φορτίο, tuned dropout/lr/wd
\end{itemize}


\section{Στρατηγικές Πρόβλεψης}
\label{sec:strategies}

\subsection{Teacher-Forcing (TF / Closed-Loop)}

Το μοντέλο εκπαιδεύεται για πρόβλεψη ενός βήματος:
\begin{equation}
    \hat{y}_{t+1} = f\bigl(\mathbf{x}_t,\, y_t,\, y_{t-1},\, y_{t-23},\, y_{t-47},\, y_{t-167}\bigr)
\end{equation}
όπου $\mathbf{x}_t$ περιλαμβάνει τα εξωγενή χαρακτηριστικά. Κατά την αξιολόγηση χρησιμοποιούνται
πραγματικές καθυστερημένες τιμές, επομένως \textit{δεν υπάρχει συσσώρευση σφάλματος}.

\subsection{Recursive (Open-Loop) με Scheduled Sampling}

Στη Recursive στρατηγική, για $h > 1$:
\begin{equation}
    \hat{y}_{t+h} = f\bigl(\mathbf{x}_{t+h-1},\, \hat{y}_{t+h-1},\, \hat{y}_{t+h-2},\, \ldots\bigr)
\end{equation}
Η αξιολόγηση γίνεται σε daily chunks (24 ώρες) ή weekly chunks (168 ώρες).

Η τεχνική \textbf{Scheduled Sampling}~\cite{bengio2015scheduled} εκπαιδεύει το μοντέλο με
πιθανότητα $\varepsilon_k$ να χρησιμοποιήσει τη δική του πρόβλεψη αντί της πραγματικής τιμής:
\begin{equation}
    \varepsilon_k = \min\!\left(\varepsilon_{\text{final}},\; \varepsilon_{\text{start}} +
    \frac{k-1}{K-1}\,(\varepsilon_{\text{final}} - \varepsilon_{\text{start}})\right)
\end{equation}
με $\varepsilon_{\text{start}}=0.10$, $\varepsilon_{\text{final}}=0.40$ σε $K=3$ iterations. Αυτό
βελτιώνει την ανθεκτικότητα του μοντέλου σε θόρυβο εισόδου κατά την recursive αξιολόγηση.

\subsection{Multi-Input Multi-Output (MIMO)}

Ένα μοντέλο εκπαιδεύεται να παράγει ταυτόχρονα $H=24$ τιμές:
\begin{equation}
    [\hat{y}_{t+1}, \hat{y}_{t+2}, \ldots, \hat{y}_{t+24}] = f\bigl(\mathbf{x}_t,\,
    y_t,\, y_{t-23},\, y_{t-47},\, y_{t-167}\bigr)
\end{equation}
Για τα δέντρα (RF, XGB, LGBM) χρησιμοποιείται \texttt{MultiOutputRegressor}.

\textbf{Two-Stage MIMO για τιμή:} Αρχικά εκπαιδεύεται MIMO μοντέλο φορτίου ($H=24$), στη
συνέχεια τα 24 forecast φορτίου ενσωματώνονται ως επιπλέον χαρακτηριστικά (176 ή 194 features)
για την εκπαίδευση MIMO μοντέλου τιμής. Αυτό αντικατοπτρίζει την αιτιώδη σχέση φορτίου--τιμής.

\subsection{Walk-Forward Monthly Retraining}

\begin{equation}
    \mathcal{D}_m = \{(x_t, y_t) : t \leq \text{end of month}\;(m-1)\}
\end{equation}
Κάθε μήνα $m$, το μοντέλο επανεκπαιδεύεται στο σύνολο $\mathcal{D}_m$ και αξιολογείται στον
μήνα $m$. Η εκπαίδευση χρησιμοποιεί κυλιόμενο παράθυρο ή expanding window (επέκταση από Ιαν.\ 2017).

Για την Teacher-Forcing παραλλαγή (TF-MR), χρησιμοποιούνται πραγματικές lagged τιμές.
Για τη Recursive παραλλαγή (Rec-SS-DO-MR), συνδυάζεται Scheduled Sampling και Daily Optuna
(Optuna ανά μήνα, 15 trials, 2 splits).


\section{Βελτιστοποίηση Υπερπαραμέτρων (Optuna)}
\label{sec:optuna}

Χρησιμοποιείται το πλαίσιο Optuna~\cite{optuna2019} για αυτόματη βελτιστοποίηση υπερπαραμέτρων
με Tree-structured Parzen Estimator (TPE) sampler.

\subsection{OL-Aware Daily Optuna}

Η συνάρτηση αξιολόγησης (objective function) μιμείται την πραγματική recursive αξιολόγηση:
κάθε fold αξιολογεί σε ωριαία chunks 24 ωρών χρησιμοποιώντας τις προβλέψεις ως lagged εισόδους.
Αυτό αποτρέπει \textit{CV leakage}: το CV χρησιμοποιεί recursive predictions (όχι teacher-forcing),
οπότε η βελτίωση CV αντανακλά πράγματι βελτίωση στο test.

Αντίθετα, το standard Optuna με TF-CV (χρήση πραγματικών τιμών) \textit{παρουσιάζει CV leakage}
για OL μοντέλα: βελτίωση CV δεν μεταφέρεται σε βελτίωση test (παρατηρήθηκε για LGBM Dense OL).

\subsection{MIMO Optuna}

Για MIMO, χρησιμοποιείται MultiOutputRegressor και $k$-fold CV με rolling-origin splits,
αξιολογώντας τον μέσο MAE για όλα τα 24 βήματα ορίζοντα.

\subsection{CachedOptuna}

Για Monthly Retraining, οι υπερπαράμετροι βελτιστοποιούνται \textit{μία φορά} στο πλήρες
training set (train\_end=Νοε.\ 2025, 40--60 trials) και αποθηκεύονται ως JSON. Κάθε μηνιαία
επανεκπαίδευση χρησιμοποιεί τις αποθηκευμένες παραμέτρους χωρίς νέο Optuna. Αυτό επιταχύνει
σημαντικά τη διαδικασία και αποφεύγει overfitting του Optuna σε μικρά validation sets.


\section{Μετρικές Αξιολόγησης}
\label{sec:metrics}

Χρησιμοποιούνται τρεις μετρικές:

\begin{description}
    \item[\textbf{MAE (Mean Absolute Error):}]
    \begin{equation}
        \text{MAE} = \frac{1}{T} \sum_{t=1}^{T} |y_t - \hat{y}_t|
    \end{equation}
    Ερμηνεύεται άμεσα σε μονάδες τιμής (€/MWh) ή φορτίου (MW).

    \item[\textbf{RMSE (Root Mean Squared Error):}]
    \begin{equation}
        \text{RMSE} = \sqrt{\frac{1}{T} \sum_{t=1}^{T} (y_t - \hat{y}_t)^2}
    \end{equation}
    Τιμωρεί πιο αυστηρά τα μεγάλα σφάλματα (outliers).

    \item[\textbf{sMAPE (symmetric Mean Absolute Percentage Error):}]
    \begin{equation}
        \text{sMAPE} = \frac{100\%}{T} \sum_{t=1}^{T}
        \frac{|y_t - \hat{y}_t|}{\frac{1}{2}(|y_t| + |\hat{y}_t|)}
    \end{equation}
    Παρέχει κανονικοποιημένο σφάλμα (ποσοστό), χρήσιμο για σύγκριση μεταξύ περιόδων
    διαφορετικής μέσης τιμής.
\end{description}

Ο Weron (2014)~\cite{weron2014electricity} και ο Lago et al.\ (2021)~\cite{lago2021} συνιστούν τη
χρήση MAE ως κύριας μετρικής για EPF, καθώς είναι ρωμαλέα σε ακραίες τιμές και αποφεύγει τα
προβλήματα διαίρεσης με το μηδέν που παρουσιάζει το MAPE. Η παρούσα εργασία ακολουθεί αυτή τη
σύσταση.
